% Section 5 Header
\section{Pattern 3: Monotonic Stack}

\begin{frame}[c]{\empty}
    \begin{center}
        \Large\textbf{\secname}
    \end{center}
\end{frame}

% Slide 1: Concept
\begin{frame}{What is a Monotonic Stack?}
  \begin{itemize}
    \item \textbf{The Core Idea:} A stack that maintains elements in a strictly increasing or decreasing order.
    \item \textbf{How it works:} When a new element is processed, we pop all elements from the stack that violate the monotonic property before pushing the new element.
  \end{itemize}
\end{frame}

\begin{frame}{Monotonic Stack: Types}
  \begin{itemize}
    \item \textbf{Monotonic Increasing Stack:} Stack elements are sorted in ascending order from bottom to top (e.g. `[1, 3, 5, 8]`). Useful to find the first smaller element to the left/right.
    \item \textbf{Monotonic Decreasing Stack:} Stack elements are sorted in descending order from bottom to top (e.g. `[9, 7, 4, 2]`). Useful to find the first larger element to the left/right.
  \end{itemize}
\end{frame}

% Problem 1: Daily Temperatures
\begin{frame}[fragile]{Problem: Days to Warmer Temperature}
  \textbf{Scenario:}
  Given an array of daily temperatures, calculate the number of days you must wait until a warmer day occurs. If no warmer day is possible, output 0.
  
  \vspace{0.3cm}
  \textbf{Examples:}
  \begin{itemize}
    \item \textbf{Input:} \texttt{[73, 74, 75, 71, 69, 72, 76, 73]} $\rightarrow$ \textbf{Output:} \texttt{[1, 1, 4, 2, 1, 1, 0, 0]}
  \end{itemize}
\end{frame}

% Problem 2: Next Greater Element I
\begin{frame}[fragile]{Problem: Next Greater Element Search}
  \textbf{Scenario:}
  Given two arrays nums1 and nums2 where nums1 is a subset of nums2, find the first element in nums2 to the right of each nums1 element that is strictly greater.
  
  \vspace{0.3cm}
  \textbf{Examples:}
  \begin{itemize}
    \item \textbf{Input:} \texttt{array1 = [4,1,2], array2 = [1,3,4,2]} $\rightarrow$ \textbf{Output:} \texttt{[-1, 3, -1]}
  \end{itemize}
\end{frame}

% Problem 3: Trapping Rain Water
\begin{frame}[fragile]{Problem: Volume of Trapped Water}
  \textbf{Scenario:}
  Given an elevation map of bar heights where the width of each bar is 1, calculate the total volume of water trapped between the bars after a heavy rain.
  
  \vspace{0.3cm}
  \textbf{Examples:}
  \begin{itemize}
    \item \textbf{Input:} \texttt{[0, 1, 0, 2, 1, 0, 1, 3, 2, 1, 2, 1]} $\rightarrow$ \textbf{Output:} \texttt{6}
  \end{itemize}
\end{frame}

% Problem 4: Largest Rectangle in Histogram
\begin{frame}[fragile]{Problem: Largest Rectangular Area in Histogram}
  \textbf{Scenario:}
  Given an array of integers representing heights of consecutive bars in a histogram, find the area of the largest rectangle that can fit inside the histogram boundary.
  
  \vspace{0.3cm}
  \textbf{Examples:}
  \begin{itemize}
    \item \textbf{Input:} \texttt{[2, 1, 5, 6, 2, 3]} $\rightarrow$ \textbf{Output:} \texttt{10}
  \end{itemize}
\end{frame}

% Common Triggers Slide
\begin{frame}{\secname\ -- Common Triggers}
  \begin{itemize}
    \item Finding the "next greater" or "previous smaller" element in an array.
        \item Computing areas bounded by heights (e.g. histograms, trapping water).
        \item Resolving queries in $O(N)$ total time where each index is pushed/popped at most once.
  \end{itemize}
\end{frame}
